\documentclass[12pt,a4paper]{article}

\usepackage[margin=2.5cm]{geometry}

\title{Volatility Forecasting: A Statistical and Deep Learning Approach}
\author{}
\date{}

\begin{document}

\maketitle

\begin{abstract}
% Maximum one page. Write this section last.
\end{abstract}

\section{Introduction}
% Write this section last.

\subsection{Context}
% Briefly describe the context of the work.
% Example: All human beings need food to survive.

\subsection{Motivations}
% Explain the problem that motivates the work.
% Example: Unfortunately, there are regions on Earth where there is a shortage of food.

\subsection{Contributions}
% Describe the contributions that address the gaps identified in the motivations.
% Example: We present a method to generate food from thin air.

\subsection{Related Work}
% Describe the state of the art, including relevant algorithms and tools.
% For every paper or tool, explain what this work provides that is not already available.
% When possible, list similar existing tools and explain how this work differs from them.

\subsection{Outline}
% Give an outline of the thesis structure using one short sentence per section.

\section{Background}
% Present all background knowledge required to understand the work.

\section{Methods}
% Describe the system design and methodology. Do not include source code here.

\section{Implementation}
% Describe how the system was implemented.
% Use pseudocode and/or informal or UML-style diagrams where useful.

\section{Experimental Results}
% Describe the experimental results. This is the most important section.

\subsection{Goals}
% Describe the objectives of the experimental activity, typically:
% - demonstrate correctness;
% - evaluate computational performance, such as CPU/GPU time and RAM usage;
% - evaluate effectiveness, such as expected gain and its standard deviation.

\subsection{Setting}
% Describe the experimental setting: hardware, software, datasets, splits,
% hyperparameters, reproducibility controls, and evaluation metrics.

\subsection{Case Studies}
% Describe the case studies used in the experiments.

\subsection{Correctness}
% Describe experiments that demonstrate the correctness of the implementation.

\subsection{Computational Evaluation}
% Evaluate computational performance, such as training time and CPU, GPU,
% RAM, and VRAM usage.

\subsection{Technical Evaluation}
% Evaluate the effectiveness of the proposed approach.
% Describe the experiments and summarize the results using tables and graphs.

\section{Conclusions}
% Summarize the work and outline possible future developments.

\begin{thebibliography}{99}
% Add one \bibitem entry for every paper or book cited in the main text.
\end{thebibliography}

\end{document}
